% ===================================================================
%  Αρχείο: 37823_SOLUTION.tex (Fragment)
%  Αυτόματη μετατροπή μέσω doc2latex.py
% ===================================================================

ΘΕΜΑ 4

Δίνεται οξυγώνιο ισοσκελές τρίγωνο ΑΒΓ (ΑΓ=ΒΓ). Η μεσοκάθετη \emph{ε} της ΑΓ τέμνει την προέκταση της ΑΒ (προς το μέρος του Β) στο σημείο Μ και την ΑΓ στο Ζ. Στην προέκταση της ΜΓ (προς το μέρος του Γ) παίρνουμε σημείο Ε τέτοιο ώστε ΓΕ=ΒΜ.

\begin{alist}
\item Να δείξετε ότι το τρίγωνο ΑΜΓ είναι ισοσκελές. \monades{8}

\item Να δειχτεί ότι τα τρίγωνα ΑΓΕ και ΓΒΜ είναι ίσα. \monades{10}

\item Να δειχτεί ότι το τρίγωνο ΑΜΕ είναι ισοσκελές. \monades{7}


% TODO: Μετατροπή σχήματος σε TikZ/PGFPlots
\begin{figure}[htbp]
\centering
\includegraphics[width=0.8\linewidth]{../extracted/37823_SOLUTION/image1.png}
% \caption{Σχήμα 1}
\end{figure}
\end{alist}
